\chapter{The bug museum}
\label{app:trouble}

Every trap in this book, indexed by \textbf{what you actually see}. These are
real bugs from the real build, not hypotheticals.

\section{How to debug this stack at all}

Before the tables --- the method, which matters more than any individual fix.

\subsection{Bisect by layer}

The stack has six layers and two seams. Find which layer is lying:

\begin{center}
\begin{tikzpicture}[font=\footnotesize\sffamily, node distance=3.5mm,
   q/.style={draw=brewergrey, rounded corners=2pt, minimum width=68mm,
             minimum height=7.5mm, inner sep=3pt, anchor=west}]
  \node[q] (a) at (0,0)     {\code{version() == 0x92}? \hfill \itshape SPI + driver (Ch.~2--3)};
  \node[q] (b) at (0,-9mm)  {\code{isCardPresent()} true? \hfill \itshape transceive + init (Ch.~4)};
  \node[q] (c) at (0,-18mm) {\code{activate()} returns a UID? \hfill \itshape anticollision + CRC (Ch.~6)};
  \node[q] (d) at (0,-27mm) {GetVersion returns \hex{91AF}? \hfill \itshape ISO-DEP framing (Ch.~7)};
  \node[q] (e) at (0,-36mm) {vectors green? \hfill \itshape cryptography (Ch.~8--11)};
  \foreach \x/\y in {a/b, b/c, c/d, d/e}{
     \draw[-{Stealth[length=3.5pt]}, brewergrey] (\x.south west) ++(4mm,0) -- ++(0,-1.5mm);}
  \node[font=\scriptsize\itshape, color=brewerred, anchor=west] at (72mm,-18mm)
       {the first \emph{no}};
  \node[font=\scriptsize\itshape, color=brewerred, anchor=west] at (72mm,-23mm)
       {is your bug};
\end{tikzpicture}
\end{center}

\textbf{Do not debug above a layer that has not answered ``yes''.} Most long
debugging sessions are someone investigating layer five while layer two is
broken.

\subsection{Three more rules}

\begin{description}[leftmargin=0pt, style=nextline]
  \item[Crypto never gets debugged on hardware.] If a card operation fails, run
        the offline vectors first. Green vectors mean the bug is in framing,
        state, or the card's configuration --- not in your mathematics. That
        single distinction saves hours.
  \item[Print the MAC \emph{input}, not the MAC.] When a MAC mismatches, the
        output tells you nothing --- it is noise either way. The input tells you
        everything: you can see immediately whether the counter, the TI, or the
        data is wrong.
  \item[Suspect the counter first.] In secure messaging, roughly nine failures
        out of ten are the command counter being off by one, not the
        cryptography.
\end{description}

\section{Hardware and SPI (Chapters~\ref{ch:spi}--\ref{ch:driver})}

\begin{traptable}
\code{version()} $= -110$ & Signed byte. Add \code{\& 0xFF}. \\
\code{version()} $=$ \hex{0x00} & MISO stuck low: power (3.3\,V!), wiring, CS pin. \\
\code{version()} $=$ \hex{0xFF} & MISO floating: not connected. \\
Version flickers between runs & Wire length / noise. \textbf{Shorten the jumpers.} Not a code bug. \\
Reads fine, writes ignored & Read bit \hex{0x80} set on writes. \\
\code{NoClassDefFoundError: com/sun/jna} & JNA missing. \code{mvn dependency:copy-dependencies}. \\
\code{No gpio chip found} in a \emph{unit test} & The test constructs real hardware --- you inherited instead of composing. \\
Second run fails, first worked & Device file still locked. \code{close()} / try-with-resources. \\
\code{Permission denied} on \code{/dev/spidev0.0} & \code{sudo usermod -aG spi,gpio \$USER}. \\
\end{traptable}

\begin{trap}{The \texttt{extends Pi4jSpiLink} disaster}
Writing \code{class Rc522Driver extends Pi4jSpiLink} compiles and looks
convenient. Then every unit test constructs a Pi4J context, hunts for a GPIO
chip, and dies on your laptop.

You have welded the driver to the hardware and destroyed the seam.
\textbf{A driver HAS-A link; it is not a kind of link.}
\end{trap}

\section{Card detection and activation (Chapters~\ref{ch:reader}, \ref{ch:activation})}

\begin{traptable}
Program runs for \emph{minutes} & Ported the C's iteration counter. Use a \code{nanoTime} deadline. \\
\code{0/15} polls, version fine & Init incomplete: \reg{TxASKReg}$=$\hex{0x40}, \reg{ModeReg}$=$\hex{0x3D}. \\
Detects once, then blind forever & Card sits in HALT. Send WUPA (\hex{0x52}) after REQA. \\
Detection about half the time & Card position or wires. Flat on the antenna, short wires. \\
7-byte UID, first byte wrong & Kept the \hex{0x88} cascade tag. Use \code{ac1[1..3]}. \\
\code{select()} returns $-1$ & CRC wrong. \reg{ModeReg}$=$\hex{0x3D} (preset \hex{0x6363}). \\
\hex{SAK = 0x08} & It is a MIFARE Classic. Wrong card. \\
RATS empty & Not 14443-4, or RATS sent twice. \\
Works, then fails on retry & Card left ACTIVE. \code{softReset()} $+$ \code{antennaOn()} per tap. \\
\code{NegativeArraySizeException} & FIFO level read too early. Check the IRQ wait loop. \\
\end{traptable}

\section{ISO-DEP and APDUs (Chapter~\ref{ch:apdu})}

\begin{traptable}
Empty reply & Block number desynced. Toggle once per exchange. \\
Works exactly once per power-up & Never toggling. \code{blockNumber \^{}= 1}. \\
Reply one byte too long & PCB not stripped. \code{copyOfRange(reply, 1, ...)}. \\
\hex{6E 00} / \hex{6D 00} & ISO-7816 error. Check CLA \hex{0x90} and the trailing \hex{0x00}. \\
\hex{91 AE} on a command that used to work & Card mid-chain from \hex{91 AF}. Drain it or reset the field. \\
GetVersion byte 3 is \hex{0x12} & It is an EV1 --- no \code{AuthenticateEV2First}. \\
\end{traptable}

\section{Cryptography (Chapters~\ref{ch:cmac}--\ref{ch:secmsg})}

\begin{traptable}
CMAC empty-message vector fails only & Missing the \code{n=1, lastComplete=false} case. \\
CMAC partial-block vectors fail & Padding: \hex{0x80} then zeros, not all zeros. \\
CMAC full-block vectors fail & $K_1$/$K_2$ swapped. Full $= K_1$, padded $= K_2$. \\
\emph{All} CMAC vectors fail plausibly & \code{leftShift} carry direction. Walk \textbf{backwards}. \\
Both session keys wrong & SV layout. Print the SV and compare to the spec. \\
Session keys swapped & Labels reversed. \hex{0xA5} $=$ enc, \hex{0x5A} $=$ mac. \\
\hex{91 7E} on auth step 1 & Sent only \code{keyNo}. Must be \code{\{ keyNo, 0x00 \}}. \\
Step-2 ciphertext $\neq$ vector & Two separate CBC calls. It is one \textbf{32-byte} run. \\
Auth ``succeeds'' always & You never compared RndA$'$. Actually throw on mismatch. \\
Fake test decrypts garbage & Replies lack \hex{91 xx}. Append the status word. \\
\textbf{Auth OK, first \code{getValue} MAC fails} & \textbf{Session built without TI.} \\
Command MACs right, response MACs wrong & Increment ordering: increment \emph{before} verifying. \\
First command OK, second fails & Counter not incremented, or incremented twice. \\
Balance unchanged after debit & No commit. Working as designed. \\
\hex{91 BE} & Boundary error: debit breaches the file's lower limit. \\
\hex{91 9D} & Permission denied: wrong key for that file's access rights. \\
\hex{91 F0} & File not found: card not provisioned. \\
Everything fails after one bad run & Card mid-transaction. Power-cycle the field, re-authenticate. \\
\end{traptable}

\subsection{The three crypto bugs that cost the most time}

\begin{description}[leftmargin=0pt, style=nextline]
  \item[1. The TI one-liner.] \code{authenticateEv2First} builds the session; if
        you pass only \code{(key, rndA, rndB)}, the TI defaults to zeros.
        Authentication still \emph{succeeds} --- the TI is not used until secure
        messaging --- so the failure appears \emph{one command later}, in code
        that is completely correct.
  \item[2. Counter increment ordering.] The counter ticks after the card accepts
        the command and \emph{before} you verify the response. Backwards, and
        every command MAC is right while every response check fails.
  \item[3. Fake replies without status words.] Your card code strips two bytes
        off every reply. A fake returning bare ciphertext loses two real bytes,
        and the decrypt produces garbage that looks exactly like a crypto bug.
\end{description}

\section{Tests and tooling}

\begin{traptable}
Test passes but proves nothing & The method is never called. Assert on a \emph{return value}. \\
Test passes, hardware fails & Fake too permissive. Make it assert what it received. \\
Changed code, behaviour identical & Stale \code{.class} files. \code{mvn clean compile}. \\
Passes alone, fails in the suite & Shared mutable state. Fresh session per \code{@Test}. \\
Debit and getValue tests conflict & Same fake, different counters. Separate \code{@Test} methods. \\
\end{traptable}

\begin{trap}{The green test that tests nothing}
The most dangerous bug in this book, because it \emph{rewards} you. You add a
\code{case 0x6C} to your fake, run the suite, see green --- and never notice the
test body does not call \code{getValue}.

Guard against it deliberately. After writing a test, ask: \emph{what would I have
to break to make this fail?} If the answer is ``nothing'', it is not a test.
\end{trap}

\section{UI and deployment (Chapters~\ref{ch:kiosk}, \ref{ch:deploy})}

\begin{traptable}
UI freezes on tap & Card work on the FX thread. Use the background worker. \\
Intermittent rendering corruption & Node touched off-thread. \code{Platform.runLater} everywhere. \\
Animations crawl & Software rendering. \code{-Dprism.forceGPU=true}. \\
Touch inverted & Display rotated, input not. Rotate the input device too. \\
A tile stuck highlighted & Xwayland cursor sitting on it. \code{Cursor.NONE}. \\
Second tap always declines & Field not power-cycled per transaction. \\
Kiosk dies overnight & Unhandled exception. Broad catch $+$ \code{Restart=always}. \\
Simulator used on the Pi & Constructor threw. Read the logged message. \\
\end{traptable}

\section{Environment oddities from the real build}

Not bugs in your code, but they will cost you time anyway.

\begin{itemize}
  \item \textbf{The Pi's IP changes constantly} on DHCP. Set a static lease or
        use \code{raspberrypi.local}.
  \item \textbf{SSH host-key conflicts} after a DHCP reshuffle: \emph{password
        authentication is disabled to avoid man-in-the-middle attacks}. Fix the
        known-hosts entry rather than disabling checking permanently.
  \item \textbf{sshd dies under memory pressure} on a Zero 2 W with a JavaFX app
        running. Add swap, or do not run the kiosk while developing.
  \item \textbf{\code{git commit} hangs for minutes} on a folder synced by
        iCloud or Dropbox --- the index refresh stalls on the network
        filesystem. Keep repositories out of synced folders.
  \item \textbf{\code{timeout} is not on macOS} by default. Use \code{gtimeout}
        from coreutils, or a background job.
\end{itemize}
